\documentclass[10pt,a4paper]{amsart}
\usepackage[margin=19mm]{geometry}
\usepackage{amsmath,amssymb,mathtools}
\usepackage[scr=rsfs,cal=euler]{mathalfa}
\usepackage{xfrac}
\usepackage[unicode,hidelinks,bookmarks=false]{hyperref}
\newcommand{\ie}{i.e.\ }
\newcommand{\cf}{cf.\ }
\newcommand{\resp}{respectively\ }
\newcommand{\Subsection}{\subsection}
\providecommand{\nobreakdash}{\nobreak\hbox{-}}
\setlength{\emergencystretch}{2em}
\multlinegap0pt
\usepackage{bm}
\usepackage{bbm}
\usepackage{dsfont}
\usepackage{xspace}
\usepackage{cancel}
\usepackage{mathtools}
\usepackage{amsthm}
\makeatletter
\@ifundefined{theorem}{\newtheorem{theorem}{Theorem}}{}
\@ifundefined{lemma}{\newtheorem{lemma}[theorem]{Lemma}}{}
\makeatother
\newcommand{\D}{\mathscr{D}}
\newcommand{\R}{\mathbb{R}}
\newcommand{\dist}[1]{\operatorname{dist}\left(#1\right) }
\newcommand{\one}{\mathds{1}}
\newcommand{\supp}[1]{\operatorname{supp}\left(#1\right) }
\title[Sparse domination via the Calder\'on-Zygmund decomposition: ]{Sparse domination via the Calder\'on-Zygmund decomposition: the example of Dini-smooth kernels}
\author{Fernando Ballesta-Yag\"ue, Jos\'e M. Conde-Alonso}
\date{Source: arXiv:2509.07659v1 (CC BY 4.0)}
\begin{document}
\maketitle

\noindent\emph{Source and licence.} Sparse domination via the Calder\'on-Zygmund decomposition: the example of Dini-smooth kernels. Version arXiv:2509.07659v1 (CC BY 4.0), distributed under the Creative Commons Attribution 4.0 International licence, as recorded in the arXiv licence metadata for this exact version and shown on the abstract page https://arxiv.org/abs/2509.07659v1. Full rights notice: licenses/arxiv-2509.07659-CC-BY-4.0.txt.\par\smallskip

\noindent\textit{Editorial scope.} Counted unit: the Lemma whose source label is \texttt{CZdecompProperties} in the article (source0.tex lines 196--212, source label \texttt{CZdecompProperties}), delivered together with the article's own complete original proof of it (source0.tex lines 214--301, one \texttt{proof} environment of 88 lines). No mathematics is written, repaired or re-derived: statement and proof are the article's own text line for line. Required context retained uncounted: Whitney (lines 162--176). Every definition, display and label of those lines is retained unchanged. Omitted from this package and not redistributed: 1--161, 177--195, 213--213, 302--1305. Nothing inside a retained excerpt is omitted or altered: every retained body is delivered exactly as the article's lines are written, so no omission receipt of any kind accompanies this package. Citations: the retained excerpts carry no citation command at all, so no bibliography entry of the article is transcribed and no reference is redistributed. Reference adaptation: none was needed and none was made; every reference inside the delivered unit resolves inside it, so no unresolved reference or citation is produced. Printed slips kept literal: no printed word, symbol, hypothesis or number of the counted statement or of its proof was altered. Layout and class: the article's own class is \texttt{amsart} with packages including graphicx, amsmath, amssymb, amsfonts, mathtools, amsthm, dsfont, hyperref, mathrsfs, anysize; the wrapper is the lane's pinned \texttt{amsart} editorial wrapper at 10pt on A4, which restates the theorem-like environments the retained text needs and adds the article's own macro definitions that the retained text needs (\texttt{\textbackslash D}, \texttt{\textbackslash R}, \texttt{\textbackslash dist}, \texttt{\textbackslash one}, \texttt{\textbackslash supp}), with extra wrapper packages (bm, bbm, dsfont, xspace, cancel, mathtools). The delivered numbering is the unit's own. Assets: no retained span contains a figure environment, a graphics-inclusion command, a PDF-page inclusion, a table or a TikZ picture; no image file is copied into this package, the package declares no asset dependency and no image file is redistributed with it. Licence: version arXiv:2509.07659v1 of this article is distributed under the Creative Commons Attribution 4.0 International licence, as recorded in the arXiv licence metadata for this exact version and shown on the abstract page https://arxiv.org/abs/2509.07659v1. A full rights notice is delivered at licenses/arxiv-2509.07659-CC-BY-4.0.txt. Characters: every retained excerpt is pure ASCII, so the delivered engine, XeLaTeX with its default Latin Modern text font, reports no missing character.
\begin{theorem}\label{Whitney}
Let $\Omega\subsetneq\R^d$ be open and fix a Whitney constant $C_{\mathcal{W}}>\sqrt{d}$. There exists a family $\mathcal{P}=\{L_j\} \subset \D$ such that
    \begin{enumerate}
        \item $\bigcup_{j}L_j=\Omega$, and the cubes $L_j$ are pairwise disjoint.
        \item $(C_{\mathcal{W}}-\sqrt{d})\cdot\ell(L_j)\leq \dist{L_j,\Omega^c}\leq 2C_{\mathcal{W}}\cdot \ell(L_j)$.
    \end{enumerate}
Moreover, if we take $C_{\mathcal{W}}>\sqrt{d}(2\Lambda_d+1)$, then
    \[
    2\Lambda_dL_j\subset \Omega  \text{ for all } j,
    \]
    and
    \[
    \text{if }\Lambda_dL_j\cap \Lambda_dL_k\neq\varnothing, \text{ then }\ell(L_j)\sim \ell(L_k).
    \]
\end{theorem}

\begin{lemma}\label{CZdecompProperties}
Let $\lambda>0$ be fixed and $0\leq f\in L_{\mathrm{loc}}^{1}(\R^d)$ be such that $\supp{f}\subset Q$. Let $\Omega \subset Q$ be open. Let $\{L_j\}$ be the Whitney decomposition of $\Omega$. If
    \begin{equation}\label{propiedadMaximalEnElAbierto}
        \mathcal{M}(f)\one_{\Omega^{c}}\lesssim \lambda,
    \end{equation}
where $\Omega^{c}=\R^d\setminus{\Omega}$, then    
\begin{equation}\label{inftyNormGoodPart}
        \|g\|_{\infty}\lesssim \lambda,%\quad i=1,2;
    \end{equation}
    \begin{equation}\label{L1NormBadPart}
        \sum_j \|b_{L_j}\|_{1}\lesssim \|f\|_{L^1(Q)},
    \end{equation}
    and if $R$ is a cube such that there exists $L_{j_0}\subset R$, then
    \begin{equation}\label{1.2_Jose}
        \langle |b|\rangle_{R}\lesssim\lambda.
    \end{equation}
\end{lemma}

\begin{proof}
We start proving \eqref{inftyNormGoodPart}. Since the sets $\Omega^{c}$ and $\{L_j\}_j$ are all pairwise disjoint, we first use \eqref{propiedadMaximalEnElAbierto} to get
\[
g\one_{\Omega^c}
=
f\one_{\Omega^c}
\leq
\mathcal{M}(f)\one_{\Omega^c}
\lesssim
\lambda,
\]
On the other hand, by Theorem \ref{Whitney}, there exists $\mu\sim 1$ such that $\mu L_j\cap \Omega^{c}\neq\varnothing$. Therefore, 
    \[
    \langle f\rangle_{L_j}
    \lesssim 
    \langle f\rangle_{\mu L_j}
    \leq 
    \mathcal{M}f(x_j)\quad \text{for some }x_j\in (\mu L_j)\cap\Omega^{c}.
    \]
By \eqref{propiedadMaximalEnElAbierto}, $\mathcal{M}f(x_j)\lesssim \lambda$, so $\langle f\rangle_{L_j}\lesssim \lambda$. Therefore, 
\[
g\one_{L_j}
=
\langle f\rangle_{L_j}\one_{L_j}
\lesssim 
% \lambda\one_{L_j}
% =
\lambda.
\]
To prove \eqref{L1NormBadPart}, we just compute
\[
\|b_{L_j}\|_{L^1}
=
\|(f-\langle f\rangle_{L_j})\one_{L_j}\|_{L^1}
\leq 
\|f\one_{L_j}\|_{L^1}
+
\langle f\rangle_{L_j}\|\one_{L_j}\|_{L^1}
=
2\|f\one_{L_j}\|_{L^1}.
\]
Since the $\{L_j\}$ are pairwise disjoint and $\bigcup L_j=\Omega\subseteq Q$, we have
\[
\sum_j \|b_{L_j}\|_{1}
\leq 
\sum 2\|f\one_{L_j}\|_{L^1}
=
2\int_{\bigcup L_j} |f|
\leq
2\|f\|_{L^1(Q)}.
\]
To prove \eqref{1.2_Jose}, let $R$ be such that $L_{j_0}\subset R$ for some $j_0$. We have
\begin{align*}
    \frac{1}{|R|}\int_{R}|b(x)|\, dx
    &=
    \frac{1}{R}\int_{R}\left|\sum (f(x)-\langle f\rangle_{L_j})\one_{L_j}(x)\right|\, dx
    \\&\leq 
    % \frac{1}{|T|}\sum_j \int_{T\cap L_j} (|f(x)|+\langle f\rangle_{L_j})\, dx
    % \\&\leq 
    \frac{1}{|R|}\sum_{j}
    \int_{R\cap L_j}|f(x)|\, dx
    +
    \frac{1}{|R|}\sum_{j}\langle f\rangle_{L_j}|R\cap L_j|
\end{align*}
Since $\mu L_{j}\cap \Omega^{c}\neq\varnothing$ for every $j$ and $L_{j_0}\subset R$, then $\mu R\cap \Omega^{c}\neq \varnothing$. So if $x_i\in \mu R\cap\Omega^{c}$, then %by \eqref{propiedadMaximalEnElAbierto}
\[
\frac{1}{|R|}\sum_{j}\int_{R\cap L_j}|f|
\leq%Los L_j son disjuntos
\frac{1}{|R|}\int_{R}|f|
\lesssim 
\frac{1}{|\mu R|}\int_{\mu R}|f|
\leq 
\mathcal{M}f(x_i)
%\overset{\eqref{propiedadMaximalEnElAbierto}}{\lesssim }
\lesssim
\lambda.
\]
To bound the other sum, we just use that $\langle f\rangle_{L_j}\lesssim \lambda$ and the disjointness of $\{L_j\}$:
\[
\frac{1}{|R|}\sum_{j}\langle f\rangle_{L_{j}}|R\cap L_{j}|
\lesssim 
\frac{1}{|R|}\sum_{j}\lambda|R\cap L_{j}|
\leq  %L_j son disjuntos
\lambda\frac{1}{|R|}|R|
=
\lambda.    \qedhere
\]
\end{proof}

\end{document}
